% !TEX root = document.tex
% =====================================================================
% 第1章 引言 —— 所有软件类型都必写
%
% 写什么（按顺序）：
%   1. 编写目的：这份说明书是干什么的，一两句。
%   2. 软件简介与主要功能：软件解决什么问题、和已有工具比补了什么缺口、
%      功能清单（itemize）。这一段是全文最容易被写长的地方，控制在3段内。
%   3. 技术特点：自研的地方在哪，别写"性能优越"这类空话，写具体算法或机制。
%   4. 适用对象：谁用、需要什么前置知识，可顺带说明能力边界。
%   5. 第三方依赖声明：列许可证，明确"不对其主张著作权"。法律合规必写。
%
% 软著写作哲学：
%   - 第一章不要堆技术细节和代码，细节放第3/5章。
%   - 只写软件真实有的功能。示例命令、截图、输出都要能在源码里对应上。
%   - 不确定的背景就简略，正确比丰富重要。
%   - 软件名统一用 \Name / \NameCN / \Version，不要在正文里手写死。
% =====================================================================

\section{引言}
\label{sec:introduction}

\subsection{编写目的}
\label{ssec:intro_purpose}
本说明书用于说明"\Name"（版本号\Version）的开发背景、功能组成、实现方式和使用方法。文档内容对应当前版本的实际代码，用作安装、使用和软著登记材料中的技术说明。

\subsection{软件简介与主要功能}
\label{ssec:intro_summary}
"\Name"（以下简称\Name）面向结构化文本数据的整理场景，提供从文件读取、字段规范化、列筛选到完整性统计的一体化处理。使用者不必编写脚本，通过命令行子命令或浏览器界面即可完成常见的数据清洗任务。

\Name{}把"读取—清洗—统计—导出"串成一条固定流程，将常用操作收敛为若干条命令和一个可在浏览器中使用的页面，使同一套处理逻辑在自动脚本与交互式查看之间保持一致。

软件的主要功能包括：
\begin{itemize}
    \item 多种编码的CSV文件读取与写出
    \item 列名规范化与单元格空白清理
    \item 空行过滤与指定列筛选
    \item 逐列完整性统计与报告导出
    \item 命令行批量处理与浏览器交互操作两种使用方式
\end{itemize}

各项功能的实现方式见第五章，命令行用法见第三章，界面操作见第四章。

\subsection{技术特点}
\label{ssec:technical_features}
\Name{}的技术特点在于把清洗规则实现为可复用的纯函数，命令行与网页界面调用同一套核心模块，避免两种入口产生不同的处理结果。列名规范化采用确定性的映射与去重规则，同名列不会被覆盖；编码探测按候选列表逐个尝试，避免直接以默认编码打开非UTF-8文件。完整性统计以逐列为单位，缺失数与填充率由同一份数据计算得出。

\subsection{适用对象}
\label{ssec:intro_audience}
本软件适用于需要对表格型文本做批量整理的研究人员和数据工作者。使用命令行方式需具备基本的终端操作经验；使用浏览器界面只需会上传文件。软件面向结构化数据文件，不处理图像、音频等非文本格式，也不替代数据库系统。

\subsection{第三方依赖声明}
\label{ssec:third_party}
本软件自行实现文件读写、清洗规则、统计汇总、命令行解析与网页界面等逻辑。所依赖的第三方库著作权归各自作者所有，本软件仅按相应许可证使用，不对其主张著作权：
\begin{itemize}
    \item \textbf{Python}（PSF License）：核心开发语言。
    \item \textbf{Streamlit}（Apache-2.0）：浏览器交互界面。
\end{itemize}
上述许可证信息以各项目官方发布为准；使用前请核对对应版本的许可证条款。
